\abstract{РЕФЕРАТ}

Объем работы равен \formbytotal{lastpage}{страниц}{е}{ам}{ам}. Работа содержит \formbytotal{figurecnt}{иллюстраци}{ю}{и}{й}, \formbytotal{tablecnt}{таблиц}{у}{ы}{}, \arabic{bibcount} библиографических источников и \formbytotal{числоПлакатов}{лист}{}{а}{ов} графического материала. Количество приложений – 2. Графический материал представлен в приложении А. Фрагменты исходного кода представлены в приложении Б.

Перечень ключевых слов: дорожное покрытие, дефекты дорожного покрытия, компьютерное зрение, нейронная сеть, YOLOv8, обнаружение объектов, обработка изображений, Python, PySide6, OpenCV, NumPy, Ultralytics, пользовательский интерфейс, инференс, обучение модели, датасет, детекция объектов, ограничивающая рамка, классификация, нейросетевой анализ, графический интерфейс.

Объектом разработки является программно-информационная система обнаружения дефектов дорожного покрытия на основе технологий компьютерного зрения и нейросетевого анализа изображений.

Целью выпускной квалификационной работы является создание программно-информационной системы обнаружения дефектов дорожного покрытия.

В процессе создания системы обнаружения дефектов дорожного покрытия были проведены анализ предметной области и существующих методов обнаружения объектов, анализ архитектуры нейросетевых моделей семейства YOLO, спроектирована архитектура программной системы и пользовательского интерфейса, разработан модуль обнаружения дефектов дорожного покрытия на основе модели YOLOv8n, выполнено обучение нейросетевой модели, реализованы механизмы загрузки, обработки, отображения и сохранения изображений, реализованы средства визуализации результатов обнаружения, формирования статистики и текстового отчёта.

При разработке программно-информационной системы использовались язык программирования Python, фреймворк PySide6, библиотеки OpenCV и NumPy, а также нейросетевая модель семейства YOLOv8.

\selectlanguage{english}
\abstract{ABSTRACT}
  
The volume of work is \formbytotal{lastpage}{page}{}{s}{s}. The work contains \formbytotal{figurecnt}{illustration}{}{s}{s}, \formbytotal{tablecnt}{table}{}{s}{s}, \arabic{bibcount} bibliographic sources and \formbytotal{числоПлакатов}{sheet}{}{s}{s} of graphic material. The number of applications is 2. The graphic material is presented in annex A. Fragments of the source code are presented in annex B.

List of keywords: road surface, pavement defects, computer vision, neural network, YOLOv8, object detection, image processing, Python, PySide6, OpenCV, NumPy, Ultralytics, user interface, inference, model learning, dataset, object detection, bounding box, classification, neural network analysis, graphical interface.

The object of the development is a software and information system for detecting pavement defects based on computer vision technologies and neural network image analysis.

In the process of creating a pavement defect detection system, an analysis of the subject area and existing object detection methods, an analysis of the architecture of neural network models of the YOLO family were carried out, the architecture of the software system and user interface was designed, a module for detecting pavement defects based on the YOLOv8n model was developed, a neural network model was trained, mechanisms for loading, processing, displaying and saving images were implemented, tools for visualizing detection results, generating statistics, and a text report are implemented.

When developing the software and information system, the Python programming language, the PySide6 framework, the OpenCV and NumPy libraries, as well as the YOLOv8 neural network model were used.
\selectlanguage{russian}
